% sec6_6.tex — Section 6.6: Incorporating Serial Correlation in GMM

\subsection*{The Model and Asymptotic Results}

We have now all the tools needed to introduce serial correlation in the single-equation
GMM model of Chapter~3. Recall that $\mathbf{g}_t$ in Chapter~3 (with the subscript
``$i$'' replaced by ``$t$'') is a $K$-dimensional vector defined as $\mathbf{x}_t \cdot \varepsilon_t$, the product
of the $K$-dimensional vector of instruments $\mathbf{x}_t$ and the scalar error term $\varepsilon_t$. By
Assumption~3.3 (orthogonality conditions), the mean of $\mathbf{g}_t$ is zero. We assumed
in Assumption~3.5 that $\{\mathbf{g}_t\}$ was a martingale difference sequence. Since no serial
correlation is allowed under this assumption, the matrix $\mathbf{S}$, defined to be the asymptotic
variance of $\bar{\mathbf{g}}$ ($\equiv \frac{1}{n} \sum_{t=1}^{n} \mathbf{g}_t$), was simply the variance of $\mathbf{g}_t$. We can now allow
for serial correlation in $\{\mathbf{g}_t\}$ by relaxing Assumption~3.5 as

\medskip
\noindent\textbf{Assumption 3.5$'$ (Gordin's condition restricting the degree of serial correlation):}
\emph{$\{\mathbf{g}_t\}$ satisfies Gordin's condition. Its long-run covariance matrix is nonsingular.}

\medskip
Then by the multivariate version of Proposition~6.10, $\sqrt{n}\,\bar{\mathbf{g}}$ converges to a normal
distribution. The variance of this limiting distribution (namely, $\avar(\bar{\mathbf{g}})$), denoted
$\mathbf{S}$, equals the long-run covariance matrix,
\begin{equation}
  \mathbf{S} = \sum_{j=-\infty}^{\infty} \bm{\Gamma}_j
  = \bm{\Gamma}_0 + \sum_{j=1}^{\infty} (\bm{\Gamma}_j + \bm{\Gamma}_j'),
  \label{eq:6.6.1}
\end{equation}
where $\bm{\Gamma}_j$ is the $j$-th order autocovariance matrix
\begin{equation}
  \bm{\Gamma}_j = \E(\mathbf{g}_t \mathbf{g}_{t-j}')
  \quad (j = 0, \pm 1, \pm 2, \ldots).
  \label{eq:6.6.2}
\end{equation}
(Since $\E(\mathbf{g}_t) = \mathbf{0}$ by the orthogonality conditions, there is no need to subtract the
mean from $\mathbf{g}_t$.)

That is, $\mathbf{S}$ ($\equiv \avar(\bar{\mathbf{g}})$) equals $\bm{\Gamma}_0$ under Assumption~3.5 and $\sum_{j=-\infty}^{\infty} \bm{\Gamma}_j$ under
Assumption~3.5$'$. This is the \emph{only} difference between the GMM model with and
without serial correlation. Accordingly, all the results we developed in Sections
3.5--3.7 carry over, provided that $\mathbf{S}$ is now given by \eqref{eq:6.6.1} and $\hat{\mathbf{S}}$, some consistent
estimation of $\mathbf{S}$, is redefined to take into account serial correlation. More
specifically:
\begin{itemize}
\item The GMM estimator $\hat{\bm{\delta}}(\hat{\mathbf{W}})$ remains consistent and asymptotically normal. Its
asymptotic variance is consistently estimated by (3.5.2), provided that the $\hat{\mathbf{S}}$
there is consistent for $\mathbf{S}$ defined in \eqref{eq:6.6.1}. This estimated asymptotic variance
sometimes gets the adjective heteroskedasticity and autocorrelation consistent
(\textbf{HAC}).

\item The GMM estimator achieves the minimum variance when $\plim_{n \to \infty} \hat{\mathbf{W}} = \mathbf{S}^{-1}$
where $\mathbf{S}$ is given by \eqref{eq:6.6.1}.

\item The two-step procedure described in Section~3.5 still provides an efficient GMM
estimator, provided that $\hat{\mathbf{S}}$ calculated in \emph{Step~1} is consistent for $\mathbf{S}$.

\item With $\hat{\mathbf{S}}$ thus properly defined, the expressions for the $t$, $W$, $J$, $C$, and $LR$ statistics
of Sections~3.5--3.7 remain valid; those statistics retain the same asymptotic
distributions in the presence of serial correlation.
\end{itemize}

\subsection*{Estimating S When Autocovariances Vanish after Finite Lags}

All these results assume that you have a consistent estimate, $\hat{\mathbf{S}}$, of the long-run
variance matrix $\mathbf{S}$. Obtaining such an estimate takes some intellectual effort, particularly
when autocovariances do not vanish after finite lags.

We start our quest with consistent estimation of individual autocovariances.
The natural estimator is
\begin{equation}
  \hat{\bm{\Gamma}}_j = \frac{1}{n} \sum_{t=j+1}^{n} \hat{\mathbf{g}}_t \hat{\mathbf{g}}_{t-j}'
  \quad (j = 0, 1, \ldots, n-1)
  \label{eq:6.6.3}
\end{equation}
where
\[
  \hat{\mathbf{g}}_t \equiv \mathbf{x}_t \cdot \hat{\varepsilon}_t, \quad
  \hat{\varepsilon}_t = y_t - \mathbf{z}_t' \hat{\bm{\delta}}, \quad
  \hat{\bm{\delta}} \text{ consistent for } \bm{\delta}.
\]
If we had the true value $\mathbf{g}_t$ in place of the estimated value $\hat{\mathbf{g}}_t$ in this formula, then
$\hat{\bm{\Gamma}}_j$ would be consistent for $\bm{\Gamma}_j$ under Assumptions~3.1 and~3.2, the assumptions
implying ergodic stationarity for $\{\mathbf{g}_t\}$. But we have to use the estimated series $\{\hat{\mathbf{g}}_t\}$
rather than the true series to calculate autocovariances. For this reason we need
some suitable fourth-moment condition (which we will not bother to state) for the
estimated autocovariances to be consistent. The proof of consistency is not given
here because it is very similar to the proof of Proposition~3.4.

Given the estimated autocovariances, there are two cases to consider for the
calculation of $\hat{\mathbf{S}}$. If we know \emph{a priori} that $\bm{\Gamma}_j = \mathbf{0}$ for $j > q$ where $q$ is known and
finite, then clearly $\mathbf{S}$ can be consistently estimated by
\begin{equation}
  \hat{\mathbf{S}} = \hat{\bm{\Gamma}}_0 + \sum_{j=1}^{q} (\hat{\bm{\Gamma}}_j + \hat{\bm{\Gamma}}_j')
  = \sum_{j=-q}^{q} \hat{\bm{\Gamma}}_j
  \quad (\text{recall: } \hat{\bm{\Gamma}}_{-j} = \hat{\bm{\Gamma}}_j').
  \label{eq:6.6.4}
\end{equation}
More difficult is the other case where we do not know $q$ (which may or may not
be infinite). How can we estimate the long-run variance matrix, which involves
possibly infinitely many parameters? There are two approaches.

\subsection*{Using Kernels to Estimate S}

Recently, several procedures have been proposed to estimate the long-run variance
matrix $\mathbf{S}$.\footnote{For the results to be stated in this and the next subsection to hold, we need to impose a set of technical conditions (in addition to Gordin's condition above) that further restrict the nature of serial correlation. For a statement of those conditions and a more detailed discussion of the subject treated in this and the next subsection, see Den Haan and Levin (1996b).}
A class of estimators, called the \textbf{kernel-based} (or ``nonparametric'')
\textbf{estimators}, can be expressed as a weighted average of estimated autocovariances:
\begin{equation}
  \hat{\mathbf{S}} = \sum_{j=-n+1}^{n-1} k\!\left(\frac{j}{q(n)}\right) \cdot \hat{\bm{\Gamma}}_j.
  \label{eq:6.6.5}
\end{equation}
Here, the function $k(\cdot)$, which gives weights for autocovariances, is called a \textbf{kernel}
and $q(n)$ is called the \textbf{bandwidth}. The bandwidth can depend on the sample size $n$.
The estimator \eqref{eq:6.6.4} above is a special kernel-based estimator with $q(n) = q$ and
\begin{equation}
  k(x) = \begin{cases} 1 & \text{for } |x| \leq 1, \\ 0 & \text{for } |x| > 1. \end{cases}
  \label{eq:6.6.6}
\end{equation}
This kernel is called the \textbf{truncated kernel}. In the present case of unknown $q$,
we could use the truncated kernel with a bandwidth $q(n)$ that increases with the
sample size. As $q(n)$ increases to infinity, more and more (and ultimately all) autocovariances
can be included in the calculation of $\hat{\mathbf{S}}$. This truncated kernel-based
$\hat{\mathbf{S}}$, however, is not guaranteed to be positive semidefinite in finite samples. This is
a problem because then the estimated asymptotic variance of the GMM estimator
may not be positive semidefinite.

Newey and West (1987) noted that the kernel-based estimator can be made
nonnegative definite in finite samples if the kernel is the \textbf{Bartlett kernel}:
\begin{equation}
  k(x) = \begin{cases} 1 - |x| & \text{for } |x| \leq 1, \\ 0 & \text{for } |x| > 1. \end{cases}
  \label{eq:6.6.7}
\end{equation}
The Bartlett kernel-based estimator of $\mathbf{S}$ is called (in econometrics) the \textbf{Newey--West
estimator}. For example, for $q(n) = 3$, the kernel-based estimator includes
autocovariances up to two (\emph{not} three) lags:
\begin{equation}
  \hat{\mathbf{S}} = \hat{\bm{\Gamma}}_0
  + \Bigl(\frac{2}{3}\Bigr)(\hat{\bm{\Gamma}}_1 + \hat{\bm{\Gamma}}_1')
  + \Bigl(\frac{1}{3}\Bigr)(\hat{\bm{\Gamma}}_2 + \hat{\bm{\Gamma}}_2').
  \label{eq:6.6.8}
\end{equation}

In his definitive treatment of kernel-based estimation of the long-run variance,
Andrews (1991) has examined positive semidefinite kernels, namely, the class of
kernels (including Bartlett's) that yield a positive semidefinite $\hat{\mathbf{S}}$. He shows that,
under some suitable regularity conditions, how the speed at which $\hat{\mathbf{S}}$ converges (in mean
square) to $\mathbf{S}$ depends on the kernel and $q(n)$. The most rapid possible rate is $n^{2/5}$
(i.e., $n^{2/5} \cdot (\hat{\mathbf{S}} - \mathbf{S})$ remains stochastically bounded),\footnote{A sequence of random variables $\{x_n\}$ is said to be \textbf{stochastically bounded} and written $x_n = O_p$ if for any $\varepsilon$ there exists an $M > 0$ such that $\Pr(|x_n| > M) < \varepsilon$ for all $n$. If a sequence converges in distribution to a random variable, then the sequence is stochastically bounded.}
which occurs when $q(n)$
grows at rate $n^{1/5}$ (i.e., $q(n)/n^{1/5}$ remains bounded) and the kernel is a member of
a certain subset of positive semidefinite kernels. The Bartlett kernel is not in this
subset. For Bartlett, the most rapid possible rate is $n^{1/3}$ which occurs when $q(n)$
grows at rate $n^{1/3}$. Therefore, at least on asymptotic grounds, positive semidefinite
kernels in this subset should be preferred to Bartlett. An example of such kernels
is the \textbf{quadratic spectral} (\textbf{QS}) \textbf{kernel}:
\begin{equation}
  k(x) = \frac{25}{12\pi^2 x^2}
  \left(\frac{\sin(6\pi x/5)}{6\pi x/5} - \cos(6\pi x/5)\right).
  \label{eq:6.6.9}
\end{equation}
Since $k(x) \neq 0$ for $|x| > 1$ in the QS kernel, all the estimated autocovariances $\hat{\bm{\Gamma}}_j$
($j = 0, 1, \ldots, n-1$) enter the calculation of $\hat{\mathbf{S}}$ even if $q(n) < n - 1$.

To be sure, the knowledge of the growth rate is not enough to determine the
value of the bandwidth for any particular data of finite length. Andrews (1991)
also considered a data-dependent formula for determining the bandwidth that not
only grows at the optimal rate but also minimizes some appropriate criterion (the
asymptotic mean square error). But even in this formula, some parameter values
need to be provided by the user, so the formula does not really provide automatic
data-based selection of the bandwidth.

\subsection*{VARHAC}

The other procedure for estimating $\mathbf{S}$, called the VAR Heteroskedasticity and Autocorrelation
Consistent (\textbf{VARHAC}) estimator, was proposed by Den Haan and Levin
(1996a). The idea is to fit a finite-order VAR to the $K$-dimensional series $\{\hat{\mathbf{g}}_t\}$ and
then construct the long-run covariance matrix implied by the estimated VAR. This
is not to say that $\mathbf{g}_t$ is assumed to be a finite-order VAR. In fact, in this procedure
(and also in the kernel-based procedure), $\mathbf{g}_t$ does not even have to be a linear
process, let alone a finite-order autoregression.
The VARHAC procedure consists of two steps.

\medskip
\noindent\emph{Step~1: Lag length selection for each VAR equation.} \; In this step, for each VAR
equation, you determine the lag length by the Bayesian information criterion
(BIC). Let $\hat{g}_{kt}$ be the $k$-th element of the $K$-dimensional vector $\hat{\mathbf{g}}_t$.
The $k$-th VAR equation can be written as
\begin{equation}
\begin{aligned}
  \hat{g}_{kt} &= \phi_{11}^{(k)} \hat{g}_{1,t-1} + \phi_{12}^{(k)} \hat{g}_{1,t-2} + \cdots + \phi_{1p}^{(k)} \hat{g}_{1,t-p} \\
  &\quad + \phi_{21}^{(k)} \hat{g}_{2,t-1} + \phi_{22}^{(k)} \hat{g}_{2,t-2} + \cdots + \phi_{2p}^{(k)} \hat{g}_{2,t-p} \\
  &\quad \cdots \\
  &\quad + \phi_{K1}^{(k)} \hat{g}_{K,t-1} + \phi_{K2}^{(k)} \hat{g}_{K,t-2} + \cdots + \phi_{Kp}^{(k)} \hat{g}_{K,t-p} + e_{kt} \\
  &= \bm{\phi}_1^{(k)} \hat{\mathbf{g}}_{t-1} + \cdots + \bm{\phi}_p^{(k)} \hat{\mathbf{g}}_{t-p} + e_{kt},
\end{aligned}
\label{eq:6.6.10}
\end{equation}
where
\[
  \underset{(1 \times K)}{\bm{\phi}_j^{(k)}} = (\phi_{1j}^{(k)}, \ldots, \phi_{Kj}^{(k)})
  \quad (j = 1, 2, \ldots, p).
\]
In picking the lag length by the BIC, the maximum lag length $p_{max}$ needs
to be specified. In Section~6.4, when the BIC was introduced in the context
of estimating the true order (lag length) of a finite-order autoregression,
$p_{max}$ was set to some integer known to be greater than or equal to the
true lag length. In the present case where $\mathbf{g}_t$ is not necessarily a finite-order
VAR, Den Haan and Levin (1996a, Theorem~2(c)) show that the
VARHAC estimator is consistent for $\mathbf{S}$ when $p_{max}$ grows at rate $n^{1/3}$ and
recommend setting $p_{max}$ to $[n^{1/3}]$ (the integer part of $n^{1/3}$). To emphasize
its dependence on the sample size, we denote the maximum lag length by
$p_{max}(n)$. Let $SSR_p^{(k)}$ be the sum of squared residuals (SSR) from the OLS
estimation of \eqref{eq:6.6.10} (the $k$-th equation in the VAR) for $t = p_{max}(n) +
1, p_{max}(n) + 2, \ldots, n$.\footnote{So the sample period is fixed throughout the process of picking the lag length.}
For $p = 0$, $SSR_p^{(k)} = \sum_{t=p_{max}(n)+1}^{n} (\hat{g}_{kt})^2$, the SSR
from the regression of $\hat{g}_{kt}$ on no regressors. Since there are $pK$ coefficients
in the $k$-th equation, the BIC criterion is
\begin{equation}
  \log(SSR_p^{(k)} / n) + pK \cdot \log(n) / n.
  \label{eq:6.6.11}
\end{equation}
(The $n$ in this expression could be replaced by the actual sample size $n -
p_{max}$ with no asymptotic consequences.) Let $p(k)$ be the $p$ that minimizes
the information criterion over $p = 0, 1, \ldots, p_{max}(n)$ for the $k$-th VAR
equation. Also let $\hat{\bm{\phi}}_j^{(k)}$ ($j = 1, 2, \ldots, p(k)$) be the OLS estimate of the
VAR coefficients $\bm{\phi}_j^{(k)}$ ($j = 1, 2, \ldots, p(k)$) when \eqref{eq:6.6.10} is estimated by
OLS for $p = p(k)$.

\medskip
\noindent\emph{Step~2: Calculating the implied long-run variance.} \; Let
\[
  P \equiv \max\{p(1), p(2), \ldots, p(K)\},
\]
the largest lag length in the $K$ equations. By construction, $P \leq p_{max}(n)$.
\emph{Step~1} produces an estimated VAR which can be written as
\begin{equation}
  \underset{(K \times 1)}{\hat{\mathbf{g}}_t} =
  \underset{(K \times K)}{\hat{\bm{\Phi}}_1}\;\hat{\mathbf{g}}_{t-1}
  + \cdots +
  \underset{(K \times K)}{\hat{\bm{\Phi}}_P}\;\hat{\mathbf{g}}_{t-P}
  + \underset{(K \times 1)}{\hat{\mathbf{e}}_t}
  \quad (t = p_{max}(n) + 1, p_{max}(n) + 2, \ldots, n).
  \label{eq:6.6.12}
\end{equation}
Since the lag length differs across equations, some rows of $\hat{\bm{\Phi}}_p$ ($p =
1, 2, \ldots, P$) may be zero. The rows of $\hat{\bm{\Phi}}_p$ are related to $\hat{\bm{\phi}}_j^{(k)}$ ($j = 1, 2, \ldots,
p(k)$) obtained in \emph{Step~1} as
\begin{equation}
  k\text{-th row of } \underset{(K \times K)}{\hat{\bm{\Phi}}_p} =
  \begin{cases}
    \hat{\bm{\phi}}_p^{(k)} & \text{for } 1 \leq p \leq p(k), \\
    \mathbf{0} & \text{for } p(k) < p \leq P.
  \end{cases}
  \label{eq:6.6.13}
\end{equation}
Let $\hat{\bm{\Omega}}$ be the VAR residual variance matrix:
\begin{equation}
  \underset{(K \times K)}{\hat{\bm{\Omega}}} \equiv \frac{1}{n} \sum_{t=p_{max}(n)+1}^{n} \hat{\mathbf{e}}_t \hat{\mathbf{e}}_t'.
  \label{eq:6.6.14}
\end{equation}
Recall that the long-run variance matrix of a VAR is given by \eqref{eq:6.3.16} for
$z = 1$. The estimator of $\mathbf{S}$ implied by the estimated VAR \eqref{eq:6.6.12} is then
\begin{equation}
  \hat{\mathbf{S}} = \left[\mathbf{I}_K - \sum_{p=1}^{P} \hat{\bm{\Phi}}_p\right]^{-1}
  \hat{\bm{\Omega}}
  \left[\mathbf{I}_K - \sum_{p=1}^{P} \hat{\bm{\Phi}}_p\right]^{-1\prime}.
  \label{eq:6.6.15}
\end{equation}

This VARHAC estimator is positive semidefinite by construction. Den Haan
and Levin (1996a) show that (under suitable regularity conditions) the estimator
converges to $\mathbf{S}$ at a faster rate than any positive semidefinite kernel-based estimator
for almost all autocovariance structures of $\mathbf{g}_t$ (which is not necessarily a finite-order
VAR). In particular, if $\mathbf{g}_t$ is finite-order VARMA, then the lag order chosen by the
BIC grows at a logarithmic rate and the VARHAC estimator converges at a rate
arbitrarily close to $n^{1/2}$, which is faster than the maximum rate $n^{2/5}$ achieved by
positive semidefinite kernel-based estimators.

\subsection*{Questions for Review}

\begin{enumerate}
\item When $\mathbf{x}_t = \mathbf{z}_t$, what is the efficient GMM estimator?

\item (Consequence of ignoring serial correlation) \; Consider the efficient GMM estimator
that presumes $\{\mathbf{g}_t\}$ to be serially \emph{un}correlated. If $\{\mathbf{g}_t\}$ is in fact serially
correlated, is the estimator consistent? Efficient?
\end{enumerate}
